\section{Conclusión}

La culminación de este proyecto de retrofitting demuestra que la integración de hardware de código abierto con entornos de programación industrial como LabVIEW es una solución viable y eficiente para la rehabilitación de sistemas robóticos complejos. Se logró transformar una plataforma cerrada en un sistema modular donde el flujo de información es transparente y controlado.

La robustez del sistema se fundamenta en la sinergia entre las etapas de potencia y control. El diseño de las PCBs personalizadas eliminó el ruido eléctrico y las caídas de tensión que limitaban el rendimiento original, mientras que el protocolo binario inspirado en Modbus aseguró una comunicación bidireccional fiable. El éxito del software radica en su capacidad de gestionar grandes volúmenes de datos mediante el procesamiento por lotes, eliminando la latencia y permitiendo una supervisión en tiempo real que es crítica para la seguridad del brazo.

En conclusión, se ha entregado una infraestructura tecnológica completa. Aunque el sistema opera actualmente bajo una lógica de control manual y ejecución de trayectorias preprogramadas, la precisión en la lectura de los encoders y la eficiencia del canal de escritura VISA dejan el camino trazado para la integración inmediata de modelos matemáticos de cinemática inversa. Este proyecto no solo devuelve la funcionalidad al brazo Mitsubishi, sino que lo convierte en una plataforma de investigación y desarrollo preparada para los desafíos de la robótica moderna.